% TOP_PROBLEM: {"id": 30006816, "title": "All groups are good: the four-dimensional disc-embedding conjecture", "category_id": 7, "category": {"id": 7, "name": "topology", "display_name": "Topology", "description": "Properties preserved under continuous deformations.", "slug": "topology", "order_index": 7, "created_at": "2026-07-31T15:26:25.671Z"}, "status": "open"}
% FIRST_POSTED: 2026-09-27
% ATTEMPT_STATUS: unresolved
% !TeX program = lualatex
\documentclass[11pt]{article}
\usepackage[margin=25mm]{geometry}
\usepackage{fontspec}
\setmainfont{Latin Modern Roman}
\usepackage{amsmath,amssymb,mathtools}
\usepackage{unicode-math}
\setmathfont{Latin Modern Math}
\usepackage{xurl,hyperref}
\hypersetup{colorlinks=true,urlcolor=blue}
\setcounter{secnumdepth}{0}
\setlength{\parindent}{0pt}
\setlength{\parskip}{5pt}
\setlength{\emergencystretch}{3em}
\begin{document}
\section*{\texorpdfstring{All groups are good: the four-dimensional disc-embedding conjecture}{All groups are good: the four-dimensional disc-embedding conjecture}}\label{30006816}
\subsection{Definitions and mathematical statement}
A capped grope is a surface tower built by attaching surfaces along symplectic basis curves and capping its tips by disks, with the standard four-dimensional thickening and framing. Height $1.5$ refers to the source's half-stage attachment convention. A double-point loop of an immersed disk follows paths on its two sheets through a self-intersection and is based using a path in the thickening. The selected definition calls $\Gamma$ good if every based height-$1.5$ disk-like capped grope and every homomorphism
$\phi:\pi_1(G^c)\to\Gamma$ admit an immersed disk with the prescribed framed boundary whose double-point loops all lie in $\ker\phi$. The conjecture is that every group is good. The source's equivalence to goodness of $F_2$ is retained in the quotation; its disk-embedding consequence concerns locally flat topological disks, not arbitrary smooth embeddings.
\par\smallskip\textit{Formulation and concept references:} \href{https://bpb-us-e2.wpmucdn.com/websites.umass.edu/dist/b/22144/files/2026/04/K3-problem-list-watermarked.pdf}{[S1]}.

\subsection{Short English statement}
Can every prescribed capped-grope configuration be replaced by an immersed disk whose double-point loops become trivial in the target group, for every group?
\subsection{Research attempt}
\paragraph{Scope and literature input.}
GPT-6 Astra Ultra prepared this unresolved attempt on 27 September 2026
using the 2026 primary problem list [S1], group-word calculations and
checks of the disk quantifiers. The source's Section 4.5 defines goodness
through framed immersed disks in capped-grope thickenings and records
closure properties and known subexponential-growth cases. Its definition
is stronger than algebraic cancellation of a signed intersection sum.
The argument below proves elementary reductions, reconstructs the
reduction to two generators using the specified known closure input,
and tests why finite quotients or abelianization do not finish it.
No smooth disk-embedding theorem is claimed.

\paragraph{What is required for one fixed grope.}
Fix a based finite height-$1.5$ capped grope $G^c$, its attaching region,
and a homomorphism $\phi:\pi_1(G^c)\to\Gamma$. The attaching circle
is null-homotopic in the thickening; the standard framed immersed-disk
existence in this setup is part of the source's discussion. Let
$g_1,\ldots,g_r$ be the based double-point loops of one such immersed
disk. The target is to choose a disk for which
\[
                        \phi(g_i)=1\quad(1\le i\le r).
\]
A different choice of basing paths conjugates the relevant group elements,
and exchanging sheets inverts them. Membership in $\ker\phi$ survives
both operations because that kernel is a normal subgroup. Thus the
condition is geometrically meaningful despite those choices. It is
not a request that a preferred written word have zero exponent sum.

If $\phi$ is the trivial homomorphism, every double-point loop already
passes the test. This settles that instance even when $\Gamma$ itself
is complicated. More generally, it is the actual image of the particular
homomorphism, rather than every element of the ambient target group,
that controls the application of a known goodness theorem.

\paragraph{Two direct inheritance lemmas.}
If $\Gamma$ is good and $H\le\Gamma$, then $H$ is good. Given
$\phi:\pi_1(G^c)\to H$, compose it with the injective inclusion
$H\hookrightarrow\Gamma$ and apply goodness. The resulting disk's
loop images are trivial in $\Gamma$, and injectivity makes them
trivial already in $H$. The disk, attaching region and framing are
unchanged in this argument.

There is a useful instance-wise variant. If the homomorphism factors
as
\[
 \pi_1(G^c)\xrightarrow{\psi}\Lambda
                       \xrightarrow{\eta}\Gamma
\]
with $\Lambda$ good, then goodness for $\psi$ supplies the required
disk for $\phi=\eta\psi$. Triviality before applying $\eta$ is enough;
$\eta$ need not be injective. This is a factorization criterion, not a
proof that every map to a quotient admits such a lift. In general,
relations in the domain can obstruct lifting a homomorphism through
a surjective group map.

Since the capped grope is a finite complex with a standard thickening,
its fundamental group is finitely generated. Therefore the image of
$\phi$ is a finitely generated subgroup of $\Gamma$. If every finitely
generated subgroup of $\Gamma$ is good, factor through that image
and apply the preceding criterion. This proves that it suffices to
settle all finitely generated groups, without assuming compact
four-manifolds realize arbitrary infinitely presented groups.

\paragraph{Why the free group on two generators is the universal test.}
The nontrivial established closure input from [S1] is that quotients
of good groups are good. I do not derive it merely from the
factorization lemma: an arbitrary map into a quotient has not been
lifted by that lemma. With this known closure result in hand, the
reduction to $F_2=\langle a,b\rangle$ is elementary.

For each $r$, the elements
\[
                 a^jba^{-j}\qquad(1\le j\le r)
\]
freely generate a subgroup of rank $r$. To check independence, write
a nonempty reduced word in these generators, grouping consecutive
powers of the same generator. Between powers of distinct generators,
expansion leaves a nonzero power of $a$ separating nonzero powers of
$b$. The free reduction cannot remove all those $b$ blocks. Thus no
nonempty reduced generator word becomes the identity in $F_2$.
If $F_2$ is good, the subgroup lemma makes every $F_r$ good.
Every finitely generated group is a quotient of an $F_r$, so quotient
closure makes it good. The preceding finite-image argument then handles
an arbitrary target group. The reverse implication is immediate because
$F_2$ is itself a group.

This reconstruction clarifies the content of the selected universal
statement: it is not made easier by restricting to finite presentations.
It also identifies exactly which part of the reduction rests on the
known four-dimensional goodness theory, namely quotient closure.
Free-group word calculations alone have not produced the required disks.

\paragraph{A growth collision calculation, and its geometric limitation.}
Suppose a geometric construction produced $2^h$ distinct histories at
stage $h$, each labeled by a word of length at most $ch$ in one fixed
finite generating set of a target group. Write $v(R)$ for the number
of group elements in the word ball of radius $R$. If the group has
subexponential growth, then for sufficiently large $h$,
\[
                              v(ch)<2^h.
\]
Indeed choose $\epsilon<(\log2)/c$ when $c>0$ and use
$v(R)\le e^{\epsilon R}$ for all sufficiently large $R$.
The pigeonhole principle gives two histories with the same label;
the quotient of their label words is trivial. The case $c=0$ is
immediate because the ball has one element.

In $F_2$, the positive words of length $h$ in $a,b$ are all reduced
and distinct, giving $2^h$ collision-free histories. Its full symmetric
word ball has size
\[
                      1+4\sum_{j=0}^{h-1}3^j=2\cdot3^h-1.
\]
There is consequently no corresponding collision forced by this count.
This explains a real algebraic difference between the growth regimes.
It does not prove that absence of a forced collision obstructs every
possible disk construction, or that a collision alone supplies one.
To use the estimate geometrically one must produce the histories with
the claimed length control and convert the collision into framed disk
moves without introducing uncontrolled new loops. Those requirements
are not proved here. The known subexponential result in [S1] uses
geometric work beyond this elementary counting observation.

\paragraph{Why finite quotients cannot be interchanged with one disk.}
Suppose one applies known goodness of finite groups to every finite
quotient $q:\Gamma\to Q$. The conclusion has the logical form
\[
          \text{for every }q\text{ there exists a disk }D_q
          \text{ whose loops are killed by }q\phi.
\]
The desired conclusion asks for a single disk with loops killed by
$\phi$. Residual finiteness distinguishes each nontrivial fixed loop,
but here the disk and its loop list may change with $q$. It does not
allow these two quantifiers to be exchanged.

There is an explicit word-level stress test in $F_2$. Given any finite
collection of finite quotients $q_1,\ldots,q_N$, let $m$ be the least
common multiple of the orders of $q_j(a)$. Then $a^m\ne1$ in $F_2$
but every $q_j(a^m)=1$. Thus every finite collection of quotient tests
can miss a nontrivial word that is allowed to depend on the tests.
This is not asserted to be a list realized by a particular capped-grope
disk. It refutes the bare inference that arbitrarily many finite tests,
with changing candidate data, give simultaneous triviality in the group.
One would need an additional compactness, bounded-complexity or coherent
geometric selection theorem, none of which is supplied by residual
finiteness alone.

\paragraph{Intersection sums and framing do not encode every loop condition.}
Replacing the target by its abelianization loses another essential
piece of information. The commutator $aba^{-1}b^{-1}$ is nontrivial
in $F_2$ but has zero image in its abelianization $\mathbb Z^2$.
Making all double-point loops homologically trivial would therefore
not make them trivial in the specified group. Similarly a signed
intersection expression can contain the formal cancellation
$g-g=0$ in $\mathbb Z[\Gamma]$ with $g\ne1$. Vanishing of that
sum does not say that either individual loop has trivial image.
These are algebraic countertests of proposed criteria, not constructed
counterexamples to goodness.

Geometrically, pairing intersections suggests Whitney disks, but in a
four-manifold a two-dimensional Whitney disk and a two-dimensional
surface have complementary dimensions. General position allows isolated
new intersection points; it does not make the disks disjoint from the
surfaces automatically. The boundary framing is an additional condition
which a move must preserve or repair. A higher-dimensional Whitney-trick
argument that discards these intersections and framings would miss the
specific difficulty of the four-dimensional problem. The locally flat
topological embedding consequences in the source are likewise not
statements about arbitrary smooth embedded disks.

\paragraph{Verification and outcome.}
The subgroup and finite-image reductions check the exact kernel
condition for each double-point loop. Free-word reductions and the
word-ball formula were checked finitely as illustrations of their
all-length proofs. The finite-quotient example explicitly retains a
nontrivial word and does not claim geometric realization. The first
remaining gap is a framed immersed-disk construction for every required
capped-grope map into $F_2$, with all double-point labels genuinely
trivial. Neither counting, residual finiteness nor algebraic
intersection cancellation establishes it. The all-groups-good
conjecture remains \textbf{UNRESOLVED}.

\subsection{Sources}
\begin{itemize}
\item[S1]Baykur, Kirby and Ruberman (eds.), K3 – A New Problem List in Low-Dimensional Topology, preliminary2026author version; Section4.5 scribed by Powell and Ray.
\url{https://bpb-us-e2.wpmucdn.com/websites.umass.edu/dist/b/22144/files/2026/04/K3-problem-list-watermarked.pdf}.
\textit{Formulation source.}
\end{itemize}
\end{document}
